\subsubsection{Alphabet décimal et congruence de contrôle de l'émetteur}
\label{sec:rec27-codigo}

On conserve la loi émettrice de \eqref{act21:tpk:bloque-decimal},
avec les lectures APP et l'observable \(\ell_9\) de
\eqref{act21:tpk:logaritmo}, les curseurs de
\ref{act21:tpk:operaciones} et le calendrier de six fenêtres de
\ref{act21:tpk:emision}. Les accumulateurs sont \(S_m,E_m,Z_m\), avec
\(Z_m=3\), et les chiffres de sortie sont notés \(d_{m,1},d_{m,2},d_{m,3}\).
La proposition~\ref{act21:tpk:recuperacion-firmas} restitue déjà les deux
signatures résiduelles. Leur combinaison avec les images effectives des
accumulateurs détermine l'alphabet suivant.

\begin{proposition}[Congruence de contrôle du troisième chiffre]
Chaque bloc émis satisfait
\[
 d_{m,3}=[5d_{m,2}-2d_{m,1}+1]_{10}.
\]
\end{proposition}
\begin{proof}
L'inverse de la proposition~\ref{act21:tpk:recuperacion-firmas}
donne \([E_m]_{10}=[d_{m,2}-d_{m,1}]_{10}\). En substituant
$E_m\equiv d_{m,2}-d_{m,1}\pmod{10}$ et $Z_m=3$ dans le troisième chiffre, on obtient
$3d_{m,1}+5(d_{m,2}-d_{m,1})+21\equiv5d_{m,2}-2d_{m,1}+1\pmod{10}$.
\end{proof}

\begin{proposition}[Alphabet décimal de quarante-deux blocs]
L'image de l'émetteur par fenêtre est constituée exactement des blocs
\[
\begin{array}{c|rrrrrrr}
[E_m]_{10}&\multicolumn{7}{c}{U_m}\\ \hline
0&114&227&443&556&669&885&998\\
1&129&232&458&561&674&890&903\\
3&149&252&478&581&694&810&923\\
5&169&272&498&501&614&830&943\\
7&189&292&418&521&634&850&963\\
9&109&212&438&541&654&870&983
\end{array}
\]
\end{proposition}
\begin{proof}
Une trajectoire cardinale lit trois termes consécutifs du cycle additif
\[
 (2,3,4,5,6,7,8,9,1).
\]
Leurs sommes sont
$(9,12,15,18,21,24,18,12,6)$ ; inverser la direction conserve l'ensemble
des sommes. Par conséquent
\[
 \operatorname{Im}S_m=\{6,9,12,15,18,21,24\},\qquad
 \operatorname{Im}[S_m]_{10}=\{1,2,4,5,6,8,9\}.
\]
Dans le canal multiplicatif, un facteur fixe divisible par trois apporte zéro. Pour les facteurs inversibles, les sommes cycliques de trois lectures sont
\[
\begin{array}{c|c}
\text{facteur fixe}&\text{sommes possibles}\\\hline
1,8&\{1,3,7,9\}\\
2,7&\{3,5,9\}\\
4,5&\{1,5,7\}
\end{array}
\]
Il en résulte $\operatorname{Im}E_m=\{0,1,3,5,7,9\}$.
L'atteignabilité peut déjà être vérifiée dans la première fenêtre : tous les
curseurs des deux tableaux suivants ont pour direction $E$.
\[
\begin{array}{c|rrrrrrr}
S_1&6&9&12&15&18&21&24\\\hline
(i_+,j_+)&(0,8)&(0,0)&(0,1)&(0,2)&(0,3)&(0,4)&(0,5)
\end{array}
\]
\[
\begin{array}{c|rrrrrr}
E_1&0&1&3&5&7&9\\\hline
(i_\times,j_\times)&(2,0)&(0,0)&(0,1)&(1,1)&(0,2)&(0,4)
\end{array}
\]
L'indépendance des deux curseurs réalise chacun des $7\cdot6$
couples. Les deux premiers chiffres restituent le couple résiduel, de sorte que ses
42 images sont distinctes. La congruence de contrôle fixe le troisième chiffre
et produit exactement le tableau énoncé.
\end{proof}

Le bloc décimal conserve ainsi les deux signatures résiduelles au moyen d'un codage injectif et d'un chiffre de contrôle. Ce résultat décrit l'alphabet de l'émetteur fini ; le prolongement des lectures régionales conserve en outre les conditions de phase, d'orientation et de mémoire de l'état enrichi.

L'image par fenêtre comporte donc \(42\) blocs. La concaténation
compatible de six fenêtres conserve le recensement de \(468\) émissions de
la proposition~\ref{act21:tpk:468}, dont la réduction ternaire produit
\(243\) mots dans l'espace ambiant de \(729\). Chaque cardinal compte
l'objet de son étape et conserve la provenance du même émetteur.
